\subsection{Duties to clients and to employers}

\begin{qfrule}[Applying Standard III, duties to clients: loyalty, prudence and care; fair dealing; suitability; performance presentation; and confidentiality, with conforming and violating cases]\label{rule:standard-duties-to-clients}
\qffield{Statement} \emph{III(A) Loyalty, prudence and care.} Act for the benefit of clients and place their interests before the employer's and one's own, with the care and skill a prudent person in the same role would use. Identify the client: for a pension plan or trust, the duty runs to the beneficiaries; for a fund managed to a mandate, to the mandate. Client brokerage belongs to the client: soft commissions must buy research that benefits the client, and directed brokerage needs a warning that best execution may suffer. Proxies have economic value and must be voted in an informed way, subject to cost--benefit judgement. Judge each investment in the context of the whole portfolio. The Standard does not make every member a legal fiduciary, but sets a minimum regardless of local law \citep[pp.~65--70]{standards-practice-handbook-2024}.

\emph{III(B) Fair dealing.} Treat clients fairly, not necessarily equally, in disseminating recommendations and taking investment action. Premium service levels are allowed if disclosed and offered to all. Changes in recommendations must reach clients who acted on the earlier advice. Oversubscribed new issues are pro-rated among suitable accounts, and the member forgoes allocations to themselves and their immediate family; but family-member accounts managed like other clients' accounts must not be excluded (p.~77). The test is how the account is managed, not whether the member has an interest in it: under VI(B), a family client account is treated like any other firm account even when the member is a beneficial owner, and the beneficial ownership adds the firm's preclearance and reporting requirements without changing the allocation (pp.~165, 167--168). Client consent cannot validate patently unfair allocation (pp.~76--78). Withholding a suitable IPO from a relative's fee-paying account to avoid the appearance of a conflict breaches III(A) (p.~74).

\emph{III(C) Suitability.} In an advisory relationship: inquire into experience, risk and return objectives and constraints before advising, and update regularly; judge suitability against the client's written objectives and the total portfolio. Best practice is a written IPS reviewed at least annually. For an unsolicited unsuitable trade, discuss it, document that it is against advice, and consider updating the IPS or ending the relationship. A manager of a fund to a mandate must stay within the mandate but is not responsible for each investor's suitability (pp.~86--90).

\emph{III(D) Performance presentation.} Make reasonable efforts to present performance fairly, accurately and completely; GIPS compliance is the best route (pp.~95--96).

\emph{III(E) Confidentiality.} Keep current, former and prospective clients' information confidential unless it concerns illegal activities, disclosure is required by law, or the client permits it. Best practice is a client-authorised secondary contact for vulnerable investors (pp.~100--102).

\qffield{Practice} Contrasting cases (original).
\begin{itemize}
\item \emph{Soft dollars.} A manager pays a broker higher commissions for its research platform, used for all accounts. \emph{Conforming} if the research aids client decisions and is disclosed. Using the same commissions to pay the firm's office rent \emph{violates} III(A).
\item \emph{Allocation.} A hot IPO fills 60 percent of the order. Allocating it pro rata across all accounts for which it is suitable is \emph{fair}. Giving the full amount to the three largest fee-paying clients \emph{violates} III(B). Leaving out a suitable account because it belongs to the manager's brother, a fee-paying client with his own IPS, also violates the Standards. So does cutting the fill of a fee-paying account the manager holds jointly with a parent: the ownership calls for preclearance and reporting, not a smaller share.
\item \emph{Suitability.} A 72-year-old client with an income IPS asks to put 30 percent in a leveraged crypto fund. The adviser explains the mismatch, records that the trade is against advice, and proposes updating the IPS; this \emph{conforms}. Executing silently and then citing ``client instruction'' does not.
\item \emph{Mandate.} A ``large-cap income'' fund manager buys non-dividend small caps because they look cheap. \emph{Violation} of III(C)(2).
\item \emph{Confidentiality.} An adviser suspects a client's son is draining the account of a client with dementia. If the client named the son as secondary contact, a call to someone else named in the agreement, or the procedure local law provides, is permitted. Discussing the account with a neighbour is a \emph{violation} of III(E).
\end{itemize}

\qffield{Conditions} Handbook, 12th edition. Legal fiduciary duties (for example under the US Advisers Act) may be stricter and then govern under I(A).

\qffield{Pitfalls} Treating client consent as curing an unfair allocation. Judging a single holding's risk without the rest of the portfolio. Assuming confidentiality ends with the relationship.

\qffield{See also} \hyperref[meth:preventing-client-duty-violations]{prevention}; \hyperref[meth:gips-composite-presentation]{GIPS}.

\end{qfrule}

\begin{qfmethod}[Preventing violations of Standard III: account statements and custody, best execution, written allocation and dissemination procedures, the IPS and suitability tests, performance records and confidentiality controls]\label{meth:preventing-client-duty-violations}
\qffield{Statement} \citep[pp.~70--71, 77--80, 90, 95--96, 102--103, 165]{standards-practice-handbook-2024}:
\begin{itemize}
\item \emph{Loyalty and care}: at least quarterly itemised statements to clients whose assets one controls, disclosure of where assets are held, segregation from other assets; disclosure of all forms of manager compensation; a proxy policy; best execution unless the client directs otherwise; regular reviews against governing documents.
\item \emph{Fair dealing}: recommendations to all interested, suitable clients at about the same time; few people aware before dissemination and a short interval between decision and release; a list of clients and holdings for notifying changes; documented and time-stamped orders; block trades at the same average price and commission; pro-rata partial fills with minimum lots; new issues allocated by client from advance indications; disclosure of allocation policy and of service levels.
\item \emph{Suitability}: a written IPS covering the client's type, beneficiaries, share of total assets, return objectives, risk tolerance, constraints (liquidity, cash flows, horizon, taxes, law, preferences such as ESG, proxy duties) and benchmarks; annual review; suitability tests on diversification, risk and fit with strategy.
\item \emph{Performance}: comply with GIPS or, failing that, use composites of similar portfolios, include terminated accounts, disclose simulated or gross/net results, and keep the data.
\item \emph{Confidentiality}: share only with colleagues serving the client; follow electronic-security policies; train all staff; obtain a secondary contact; extra review of vulnerable clients' accounts.
\end{itemize}

\qffield{Practice} \emph{A trade-allocation policy for Linton Partners} (original). Orders time-stamped in the order-management system; all accounts in a block receive the average price and the same commission; partial fills pro rata, rounded to 100 shares and to a USD~25,000 minimum for bonds; IPO interest gathered by client the day before pricing. Two questions are answered separately for every account, in advance.

\emph{Allocation class} (Standard III(B); VI(B) pp.~165, 167--168), decided by how the account is managed:
\begin{itemize}
\item \emph{Personal}: the employee's own accounts, and any other account in which an employee or immediate family member has a beneficial interest (ownership, control or a material financial interest) that is not managed under the client process, such as a spouse's brokerage account. They take no part in an oversubscribed issue.
\item \emph{Family client}: a relative's fee-paying account managed under the same IPS process, terms and strategy as other clients, such as a parent's retirement account, \emph{whether or not} the employee has a beneficial interest in it. It takes part in allocations exactly like other eligible accounts, neither favoured nor excluded.
\end{itemize}
\emph{Conflict controls} (VI(A), VI(B) and the firm's code), decided by the employee's interest: every family client account is flagged in the conflicts register and its fills appear in compliance's monthly report of allocation dispersion by account. If the employee also has a beneficial interest in it, the account is in addition listed in the employee's holdings report, its statements are copied to compliance, and its participation in each block is precleared before the order is placed; preclearance confirms eligibility and suitability and cannot change the pro-rata share. Trading inside the client block is not a personal trade ahead of clients, so blackout periods do not apply to it. The policy and any change are sent to clients.

Applied to the IPO case above, a 60 percent fill on orders of 5,000, 3,000 and 2,000 shares gives 3,000, 1,800 and 1,200. If the 2,000-share order belongs to the manager's father, a family client, he still receives 1,200. If instead it is a taxable investment account the manager holds jointly with her mother, managed in the same strategy on standard fees, it is a family client account with beneficial ownership: it also receives 1,200, after preclearance and with reporting. If the manager's spouse's brokerage account also asked for 1,000, it receives nothing and the clients' fills do not change.

\qffield{Conditions} Recommendations; GIPS is a voluntary standard that firms claim.

\qffield{Pitfalls} An allocation policy that is disclosed but discretionary in practice. An IPS collected once and never reviewed. A blanket exclusion of every account linked to an employee's family, which disadvantages genuine family clients. Using beneficial ownership to decide the allocation share instead of to trigger preclearance and reporting, which leaves a family client account with an employee interest unclassified.

\qffield{See also} \hyperref[rule:standard-duties-to-clients]{Standard III}; \hyperref[meth:private-client-ips]{private-client IPS}.

\end{qfmethod}

\begin{qfrule}[Applying Standard IV, duties to employers: loyalty and leaving an employer, additional compensation, and supervisors' responsibilities, with conforming and violating cases]\label{rule:standard-duties-to-employers}
\qffield{Statement} \emph{IV(A) Loyalty.} In matters related to employment, act for the employer's benefit: do not deprive it of one's skills, divulge confidential information or otherwise harm it. This is not a duty to put the employer before family and personal obligations. Independent business offering the employer's services needs the employer's consent. When leaving, preparations are allowed, but until employment ends one may not solicit clients, misuse confidential information or client lists, or take records without written permission. After leaving, clients may be contacted from public information unless an agreement forbids it. Whistle-blowing that protects clients or markets can justify acts otherwise contrary to the employer's interest \citep[pp.~107--111]{standards-practice-handbook-2024}.

\emph{IV(B) Additional compensation.} Do not accept gifts, benefits or compensation that compete with or might create a conflict with the employer's interest without written consent from all parties, including client-paid performance bonuses (pp.~117--118).

\emph{IV(C) Responsibilities of supervisors.} Make reasonable efforts to ensure that anyone under one's supervision or authority, member or not, complies with law, firm policy and the Standards. That requires an adequate compliance system, enforced and monitored. On learning of a possible violation, investigate promptly and restrict the employee meanwhile; a warning and the employee's assurances are not enough. A supervisor who cannot supervise because the system is inadequate should decline the responsibility until it is fixed (pp.~120--124).

\qffield{Practice} Contrasting cases (original).
\begin{itemize}
\item \emph{Leaving.} Nora resigns from a private bank, gives two weeks' notice, and in that period tells her clients she is leaving and names her new employer. \emph{Conforming}. Giving them her new e-mail and describing the new bank's services before her last day, or downloading the client list, \emph{violates} IV(A). After leaving, contacting former clients she finds through public directories is allowed unless her contract forbids it.
\item \emph{Moonlighting.} An analyst teaches an evening finance course for a fee. \emph{No conflict} with the employer's business, though notifying the employer is prudent. Offering paid portfolio reviews to private clients the bank could serve, without consent, \emph{violates} IV(A) and IV(B).
\item \emph{Client bonus.} A family office offers its manager a personal bonus if returns exceed 12 percent. Accepting without the employer's written consent \emph{violates} IV(B); the incentive could push the manager to take more risk in that account than in others.
\item \emph{Supervision.} A desk head learns that a trader front-ran a client order, warns him and moves on. \emph{Violation} of IV(C): she must investigate the extent, limit his trading and review the procedures that failed.
\end{itemize}

\qffield{Conditions} Applies to employees; independent contractors' duties come from their agreements (Handbook, pp.~110--111). Non-solicitation agreements and local employment law may be stricter.

\qffield{Pitfalls} Treating a resignation letter as ending the duty of loyalty. Treating the existence of a code of ethics as reasonable supervision.

\qffield{See also} \hyperref[meth:preventing-employer-duty-violations]{prevention}.

\end{qfrule}

\begin{qfmethod}[Preventing violations of Standard IV: outside-activity and departure policies, written consent for additional compensation, and a compliance system that supervisors can enforce]\label{meth:preventing-employer-duty-violations}
\qffield{Statement} \citep[pp.~111, 117--118, 122--125]{standards-practice-handbook-2024}:
\begin{itemize}
\item \emph{Loyalty}: know the employer's code and policies on outside services, termination, client notification, social media and whistle-blowing; keep personal and professional social-media accounts separate; return or erase firm records on departure.
\item \emph{Additional compensation}: disclose in writing to the supervisor or compliance the nature, approximate amount and duration of any outside compensation, and have the payer confirm the terms.
\item \emph{Supervision}: a plain-language code of ethics separate from detailed procedures; written, accessible compliance procedures that name a compliance officer with authority and resources, set the hierarchy of supervision, provide checks and balances, and document monitoring and testing; dissemination, periodic updates, training and reminders; professional conduct in performance reviews; prompt investigation and interim limits when a violation is found; and incentives that reward how results are achieved, not only revenue.
\end{itemize}

\qffield{Practice} \emph{Supervisory checklist at Marsh Capital} (original), a 25-person adviser whose new head of equities inherits a team:
\begin{enumerate}
\item Confirm the compliance manual names her as supervisor and gives her access to trade and e-mail surveillance; if not, ask senior management to fix it before accepting the role.
\item Collect annual outside-activity and compensation certifications; reject any not approved in writing.
\item Review departing employees' last month of e-mail and downloads; send the firm's client notice.
\item Replace a quarterly bonus based only on new assets with one that weights compliance findings and client complaints.
\item Log every escalated issue, the investigation and the outcome.
\end{enumerate}
Item 4 addresses the incentive that the Handbook identifies as a cause of misconduct (pp.~124--125); item 1 applies its advice to decline supervision without an adequate system (p.~121).

\qffield{Conditions} Recommendations; supervisory obligations under securities law may be specific and stricter.

\qffield{Pitfalls} Delegating supervision without telling the delegate how to detect violations. Measuring compliance by the absence of detected violations.

\qffield{See also} \hyperref[rule:standard-duties-to-employers]{Standard IV}; \hyperref[meth:asset-manager-code-compliance]{Asset Manager Code}.

\end{qfmethod}
